\begin{table}[!htbp]
\centering
\caption{Collective-reparation rollout by recorded project year}
\label{tab:desc_treatment_rollout}
\begin{tabular}{rrrr}
\toprule
Year & Newly treated & Cumulative treated & Share of RUV (\%) \\
\midrule
2007 & 260 & 260 & 4.6 \\
2008 & 353 & 613 & 10.7 \\
2009 & 488 & 1,101 & 19.3 \\
2010 & 259 & 1,360 & 23.8 \\
2011 & 165 & 1,525 & 26.7 \\
2012 & 188 & 1,713 & 30.0 \\
2013 & 117 & 1,830 & 32.0 \\
2014 & 113 & 1,943 & 34.0 \\
2015 & 93 & 2,036 & 35.6 \\
2016 & 125 & 2,161 & 37.8 \\
2017 & 248 & 2,409 & 42.2 \\
2018 & 142 & 2,551 & 44.7 \\
2019 & 533 & 3,084 & 54.0 \\
2020 & 292 & 3,376 & 59.1 \\
2021 & 121 & 3,497 & 61.2 \\
2022 & 556 & 4,053 & 71.0 \\
2023 & 168 & 4,221 & 73.9 \\
\bottomrule
\end{tabular}
\parbox{0.98\linewidth}{\footnotesize \textit{Notes:} The unit is the RUV centro poblado. A community enters its treatment cohort in the first project year recorded in the CMAN roster and remains treated thereafter. The denominator for the cumulative percentage is all 5,712 RUV communities. The 1,491 communities without a linked CMAN project by 2023 remain untreated under the current rule. Source: CMAN 2023 project roster linked to RUV Libro Segundo.}
\end{table}
